\section{Estudio técnico de operación y organización}
\label{sec:organizacion}

\subsection{Recursos, localización y capacidad}

Las capturas anteriores muestran el prototipo. Para ofrecerlo como servicio se
necesitan equipos disponibles, atención a consultas, respaldos y separación de
datos entre gimnasios. La siguiente tabla reúne esos recursos y su costo
previsto.

\begin{longtable}{p{3cm}p{3.2cm}p{4cm}p{3.7cm}}
\caption{Inventario operativo propuesto y tratamiento económico}\\
\toprule Recurso & Uso y cantidad funcional & Costo o tratamiento & Control pendiente \\
\midrule\endfirsthead
\toprule Recurso & Uso y cantidad funcional & Costo o tratamiento & Control pendiente \\
\midrule\endhead
Computadora y teléfono & Al menos un puesto de desarrollo y un dispositivo móvil de prueba. & Equipos disponibles como aportes en especie; sin compra nueva presupuestada. & Inventario físico, propietario, estado y reposición. \\
Conexión y espacio & Área de trabajo con internet y energía; operación remota inicial. & Aporte disponible; no se modela alquiler comercial. & Disponibilidad, respaldo de conexión y costo incremental. \\
Infraestructura digital & Aplicación, base de datos, dominio y servicios de ejecución. & CRC 58\,125 mensuales presupuestados. & Detallar proveedores y revisar gastos repetidos en correo y diseño. \\
Correo y diseño & Comunicación y materiales de presentación. & CRC 6\,000 mensuales presupuestados. & Revisar si Workspace ya está incluido en infraestructura. \\
Desarrollo comercial & 96 horas: aislamiento, autenticación, personalización, tareas y pruebas. & CRC 198\,735,27 de trabajo valorado. & Validar alcance, tiempo y pruebas antes de incorporar varias empresas. \\
Soporte & 0,5 horas por gimnasio al mes como supuesto base. & CRC 1\,035,08 por gimnasio de trabajo aportado. & Contrastar contra horas reales e incidencias. \\
\bottomrule
\end{longtable}
{\small\textit{Nota}. Elaboración propia a partir del escenario financiero
preliminar. Falta completar el inventario de equipos y confirmar los precios.}

Trabajar desde Pérez Zeledón facilita las visitas y las demostraciones.
Se propone comenzar de forma remota con el espacio y los equipos disponibles,
por lo que no se presupuesta alquiler. Antes de iniciar se revisará la
disponibilidad de internet, energía y equipos.

El presupuesto utiliza un tipo de cambio de CRC 500 por dólar y una referencia
horaria de CRC 496\,838,17 divididos entre 30 días y ocho horas, equivalente a
CRC 2\,070,16. Son referencias elegidas para estimar el trabajo y los servicios
en dólares. Se revisarán con precios y condiciones reales antes de contratar.

\subsection{Proceso de prestación y control de calidad}

La atención comienza con una conversación sobre las necesidades del gimnasio.
Después se realiza la demostración, se acuerda la prueba y se configura el
sistema. Durante el uso se dará seguimiento para decidir si conviene continuar.
Antes de cargar información se revisarán los permisos y se probará que los
registros se guardan y recuperan correctamente. Para atender varios gimnasios
también se comprobará que cada uno acceda solo a sus propios datos.

La compra de servicios de nube se revisará antes de contratar: finalidad,
costo, límite de uso y dependencia operativa. Se mantendrá un inventario de
cuentas y renovaciones. El almacenamiento digital se controlará por consumo,
retención y respaldo; no se trata como inventario físico de mercancías. Una
incidencia se registrará con fecha, función afectada, prioridad y responsable,
evitando almacenar contraseñas o información sensible en la bitácora.

\subsection{Estructura funcional, cargos y personal}

Derian Acuña Rojas e Ignacio Masis Cascante compartirán el desarrollo, la
atención y las decisiones del proyecto. El docente tutor, Andrés Mena Abarca,
brinda orientación académica. La tabla distribuye las tareas que asumiría el
equipo al iniciar el servicio.

\begin{table}[htbp]
\caption{Organización funcional propuesta}
\begin{tabularx}{\textwidth}{p{3.2cm}X X}
\toprule Función & Responsabilidad & Registro de seguimiento \\
\midrule
Coordinación compartida & Priorizar alcance, autorizar gastos y revisar riesgos. & Acuerdos del equipo y cierre mensual. \\
Desarrollo y calidad & Mantener aplicación, pruebas, datos y despliegue. & Repositorio, incidencias y evidencia de verificación. \\
Mercadeo y atención & Contactar, demostrar, incorporar y atender consultas. & Registro de prospectos, altas, bajas y soporte. \\
Control económico & Conciliar cobros, gastos, horas y presupuesto. & Libro financiero y comprobantes verificables. \\
Revisión académica & Orientar el informe y la presentación del proyecto. & Observaciones del tutor. \\
\bottomrule
\end{tabularx}
\caption*{\fuentepropia}
\end{table}

Las funciones técnicas y comerciales se distribuyen entre los dos integrantes;
no representan cuatro puestos adicionales. El modelo no contempla contratación
asalariada: valora el trabajo propio sin pagarlo y no lo registra como deuda.
El perfil requerido combina desarrollo web, pruebas, documentación, comunicación
con clientes y control básico del presupuesto. La necesidad de contratar
personal surgirá cuando las horas reales superen la disponibilidad acordada;
ese cambio exige incorporar remuneración y obligaciones aplicables al escenario.

\subsection{Sistema de información para dirigir el negocio}

El equipo llevará registros de contactos, pruebas, contrataciones,
cancelaciones, comprobantes y horas de trabajo. Cada mes comparará los cobros
y gastos con el presupuesto, junto con el uso de nube y el tiempo de soporte.
La información comercial tendrá acceso restringido y se mantendrá separada
de los datos de los miembros.

Cada cambio indicará quién lo registró, cuándo y con qué comprobante.
Los informes públicos usarán resultados agrupados y datos ficticios. Antes de
trabajar con información real se probará la recuperación de una copia de
seguridad.

\subsection{Forma jurídica y evidencia de formalización}

Pulso está en etapa de proyecto estudiantil y todavía no cuenta con documentos
de constitución, inscripción o afiliación. Estos trámites forman parte de los
requisitos que deben revisarse antes de operar comercialmente y de la evidencia
solicitada por la rúbrica \parencite{rubrica_plan_negocios_2026}.

Con orientación profesional, el equipo revisará si conviene operar como
persona física o constituir una empresa. También consultará cómo pueden
participar los integrantes y qué trámites tributarios, municipales y de
seguridad social corresponden. Se anotarán los costos y requisitos para
incorporarlos al presupuesto y completar la formalización.

\subsection{Alianzas y continuidad}

Además de la tutoría académica, se buscará un gimnasio dispuesto a realizar
la prueba. Su responsable aportaría tiempo y comentarios sobre las tareas;
el equipo de Pulso ofrecería configuración y soporte. El acuerdo indicaría la
duración, el uso permitido de la información y cómo terminar la prueba.

Mientras se consigue ese gimnasio, continuarán las demostraciones con datos
ficticios. El equipo revisará los motivos de rechazo para mejorar la oferta.
Los proveedores de tecnología seguirán aportando los servicios de alojamiento
y datos necesarios para la demostración.
